\chapter*{O}
\label{chap:o}

\term{Object}
An object is an entity in a category to which morphisms are connected. In the category of sets, the objects are sets; in other categories they may be groups, spaces, or vector spaces. \par\noindent\textbf{Applications:} objects provide the units that are compared and transformed in abstract mathematics.

\term{Occurrence}
An occurrence is one appearance of a symbol or term at a particular position in a formula or expression. The same symbol may have several occurrences with different meanings or scopes. \par\noindent\textbf{Applications:} occurrences are important in parsing, formal logic, compilers, and symbolic algebra.

\term{Octonion}
An octonion is an element of an eight-dimensional real division algebra that extends the complex numbers and quaternions. Multiplication is not associative, but every nonzero octonion has an inverse. \par\noindent\textbf{Applications:} octonions appear in exceptional geometry, rotations, theoretical physics, and some models of string theory.

\term{Octahedron}
An octahedron is a Platonic solid with eight equilateral triangular faces, six vertices, and twelve edges. It is dual to the cube, meaning their faces and vertices correspond. \par\noindent\textbf{Applications:} octahedral geometry appears in crystallography, molecular structure, graphics, and spatial design.

\term{Octave}
In some mathematical contexts, an octave is an informal name for an octonion. More broadly, the term can refer to a doubling step in the construction of normed algebras. \par\noindent\textbf{Applications:} this terminology occurs in discussions of composition algebras and exceptional algebraic structures.

\term{Odd Permutation}
An odd permutation is a product of an odd number of transpositions. Its sign is $-1$, whereas even permutations have sign $+1$ and form the alternating group $A_n$. \par\noindent\textbf{Applications:} permutation parity is used in determinants, sorting, combinatorics, and group theory.

\term{One-Point Compactification}
The one-point compactification adds one point at infinity to a locally compact Hausdorff space $X$. A neighborhood of infinity is the complement of a compact subset of $X$. \par\noindent\textbf{Applications:} it turns unbounded spaces into compact ones and is useful in topology, complex analysis, and dynamical systems.

\term{One-Sided Limit}
A one-sided limit describes the value approached as $x$ tends to $a$ from the right or left. A function is continuous at $a$ when both one-sided limits exist, agree, and equal $f(a)$. \par\noindent\textbf{Applications:} one-sided limits analyze jumps, boundary behavior, piecewise functions, and numerical approximations.

\term{Operad}
An operad records operations with several inputs and one output, together with rules for composing those operations. It separates the pattern of composition from the particular objects on which the operations act. \par\noindent\textbf{Applications:} operads organize homotopy theory, deformation theory, algebraic topology, and structures in mathematical physics.

\term{Open Cover}
An open cover of a space is a collection of open sets whose union contains the whole space. Compactness means that every open cover has a finite subcover. \par\noindent\textbf{Applications:} open covers define compactness, manifolds, sheaves, partitions of unity, and local-to-global arguments.

\term{Open Interval}
An open interval $(a,b)$ contains all real numbers strictly between $a$ and $b$ but excludes the endpoints. It is an open set in the usual topology on the real line. \par\noindent\textbf{Applications:} open intervals describe domains, neighborhoods, convergence ranges, and feasible strict inequalities.

\term{Open Mapping Theorem}
The open mapping theorem states that a surjective bounded linear operator between Banach spaces maps open sets to open sets. It also implies that a bijective bounded linear operator has a bounded inverse. \par\noindent\textbf{Applications:} it establishes stability and continuous solvability in functional analysis and operator equations.

\term{Open Problem}
An open problem is a precisely stated mathematical question for which no accepted proof or disproof is known. Its status can change when new research resolves it. \par\noindent\textbf{Applications:} open problems guide research and often motivate new methods, conjectures, algorithms, and theories.

\term{Open Set}
An open set contains a neighborhood of each of its points. In a metric space, every point of an open set has a small open ball contained in the set. \par\noindent\textbf{Applications:} open sets define continuity, convergence, differentiability, manifolds, and the topology of function spaces.

\term{Operation}
An operation is a rule that combines one or more inputs to produce an output, such as addition, multiplication, or composition. Its properties, such as associativity or commutativity, determine the algebraic structure. \par\noindent\textbf{Applications:} operations model combination, transformation, computation, and composition of processes.

\term{Operator (Linear)}
A linear operator is a linear map $T:V\to W$ satisfying $T(au+bv)=aT(u)+bT(v)$. On Hilbert spaces, operators may be bounded or unbounded, and their domains matter. \par\noindent\textbf{Applications:} operators represent derivatives, integrations, matrices, quantum observables, and evolution equations.

\term{Operator Algebra}
Operator algebra studies algebras of bounded operators on Hilbert spaces, including C*-algebras and von Neumann algebras. Algebraic operations and analytic norms interact in these structures. \par\noindent\textbf{Applications:} operator algebras support quantum mechanics, functional analysis, noncommutative geometry, and statistical mechanics.

\term{Operator Norm}
The operator norm is $\|T\|=\sup_{\|x\|=1}\|Tx\|$. It is the largest factor by which $T$ can stretch a unit vector and satisfies $\|Tx\|\leq\|T\|\|x\|$. \par\noindent\textbf{Applications:} it measures sensitivity, bounds errors, and analyzes stability and convergence of operators and numerical algorithms.

\term{Optimal Control}
Optimal control chooses an input $u(t)$ to minimize a cost while satisfying system dynamics and constraints. Pontryagin's principle gives necessary conditions, while dynamic programming gives Bellman or Hamilton--Jacobi--Bellman equations. \par\noindent\textbf{Applications:} optimal control is used in robotics, aerospace, energy systems, economics, and finance.

\term{Orbit (Group Action)}
The orbit of $x$ under a group action is $\operatorname{Orb}(x)=\{g\cdot x:g\in G\}$. It contains every point reachable from $x$ by the allowed symmetries. For finite groups, orbit--stabilizer gives $|\operatorname{Orb}(x)|=[G:G_x]$. \par\noindent\textbf{Applications:} orbits classify symmetry classes and count configurations in geometry, chemistry, and combinatorics.

\term{Ordinary Differential Equation (ODE)}
An ordinary differential equation relates an unknown function of one independent variable to its derivatives. Initial or boundary data select particular solutions. \par\noindent\textbf{Applications:} ODEs model motion, circuits, population dynamics, chemical reactions, and control systems.

\term{Orthogonal Complement}
For a subspace $V$ of an inner-product space $W$, the orthogonal complement is $V^\perp=\{w:\langle w,v\rangle=0\text{ for every }v\in V\}$. In finite-dimensional Euclidean spaces, $W=V\oplus V^\perp$. \par\noindent\textbf{Applications:} orthogonal complements support projections, least squares, Fourier expansions, and constraint decomposition.

\term{Orthogonal Group}
The orthogonal group $O(n)$ consists of matrices $Q$ satisfying $Q^{\mathsf T}Q=I$. Its transformations preserve inner products, lengths, angles, and distances. \par\noindent\textbf{Applications:} orthogonal groups describe rotations and reflections in geometry, mechanics, computer graphics, and physics.

\term{Orthogonal Matrix}
An orthogonal matrix satisfies $Q^{\mathsf T}Q=I$. Its columns form an orthonormal basis, its inverse is $Q^{\mathsf T}$, and its determinant is $1$ or $-1$. \par\noindent\textbf{Applications:} orthogonal matrices provide stable coordinate changes, rotations, reflections, and QR-based numerical algorithms.

\term{Orthogonal Projection}
An orthogonal projection maps a vector to its closest point in a subspace. Its matrix satisfies $P^2=P$ and $P^*=P$. \par\noindent\textbf{Applications:} projections are used in least squares, signal filtering, approximation, and separating components of data.

\term{Orthogonality}
Two vectors are orthogonal when their inner product is zero. In Euclidean space this means they are perpendicular, although the definition also works in function spaces. \par\noindent\textbf{Applications:} orthogonality simplifies coordinates, Fourier series, regression, signal separation, and numerical computation.

\term{Orthonormal Basis}
An orthonormal basis consists of mutually orthogonal unit vectors that span the space. Every finite-dimensional inner-product space has one, for example by Gram--Schmidt. \par\noindent\textbf{Applications:} orthonormal bases simplify projections, Fourier expansions, quantum states, and stable numerical representations.

\term{Osculating Circle}
The osculating circle is the circle that matches a curve's position, tangent, and curvature at a point. Its radius is $1/|\kappa|$ when the curvature $\kappa$ is nonzero. \par\noindent\textbf{Applications:} it measures local bending and is used in computer-aided design, robotics, road geometry, and motion planning.

\term{Outer Measure}
An outer measure assigns a nonnegative size to every subset, is monotone, and is countably subadditive. Carathéodory's criterion selects the sets on which the outer measure is additive, producing a measurable space. \par\noindent\textbf{Applications:} outer measures construct Lebesgue measure and handle size questions for irregular sets.

\term{Overdetermined System}
An overdetermined system has more equations than unknowns. It may have no exact solution, so least squares finds the vector that minimizes the residual $\|Ax-b\|$. \par\noindent\textbf{Applications:} overdetermined systems arise in measurement, regression, calibration, tomography, and parameter estimation.

\term{Oval}
An oval is a smooth, closed, convex plane curve. Every ellipse is an oval, but many ovals are not ellipses. \par\noindent\textbf{Applications:} ovals model rounded shapes in geometry, optics, mechanical design, and computer graphics.

\term{O-minimal Structure}
An o-minimal structure is an ordered mathematical structure in which every definable subset of the line is a finite union of points and intervals. This prevents uncontrolled one-dimensional oscillation. \par\noindent\textbf{Applications:} o-minimality supports tame geometry, real algebraic geometry, optimization, and differential equations.

\term{Order (Group Theory)}
The order of a finite group is its number of elements, written $|G|$. The order of an element $g$ is the least positive integer $n$ for which $g^n=e$, if such an integer exists. \par\noindent\textbf{Applications:} order restricts subgroup structure, periodic behavior, and symmetries in finite groups.

\term{Order (Field)}
A field order is a total order compatible with addition and multiplication: $a\leq b$ implies $a+c\leq b+c$, and nonnegative elements have nonnegative products. The real numbers admit such an order, while the complex numbers do not. \par\noindent\textbf{Applications:} ordered fields support inequalities, limits, optimization, and real analysis.

\term{Order of Growth}
Order of growth describes how a quantity behaves as its input becomes large. Big-O notation gives an upper asymptotic bound, while $\Theta$ describes matching growth. \par\noindent\textbf{Applications:} growth rates compare algorithm efficiency, resource use, convergence speed, and model scaling.

\term{Order-Preserving Map}
An order-preserving map satisfies $x\leq y\Rightarrow f(x)\leq f(y)$. It carries the order information from one partially ordered set to another. \par\noindent\textbf{Applications:} such maps represent monotone processes, lattice homomorphisms, ranking systems, and fixed-point algorithms.

\term{Order-Statistic}
The $k$th order statistic is the $k$th smallest observation after sorting a sample. The median is a central order statistic, and the minimum and maximum are the first and last. \par\noindent\textbf{Applications:} order statistics support robust summaries, quantiles, outlier detection, and nonparametric inference.

\term{Ordered Pair}
An ordered pair $(a,b)$ records both entries and their positions. Equality means $(a,b)=(c,d)$ exactly when $a=c$ and $b=d$. \par\noindent\textbf{Applications:} ordered pairs represent coordinates, relations, input-output pairs, and points in a Cartesian product.

\term{Ordering}
An ordering is a relation that organizes elements. A partial order is reflexive, antisymmetric, and transitive; a total order also compares every pair. \par\noindent\textbf{Applications:} orderings support scheduling, dependency analysis, sorting, lattices, and optimization.

\term{Ordinary Least Squares (OLS)}
Ordinary least squares estimates regression coefficients by minimizing the sum of squared residuals between observed and fitted responses. It provides coefficient estimates, fitted values, and uncertainty measures under statistical assumptions. \par\noindent\textbf{Applications:} OLS is used for prediction, calibration, trend analysis, and experimental modeling.

\term{Orthocenter}
The orthocenter is the common intersection of the three altitudes of a triangle. It lies outside an obtuse triangle and inside an acute triangle. \par\noindent\textbf{Applications:} the orthocenter is used in triangle geometry and illustrates concurrency and perpendicularity.

\term{Orthogonal Polynomials}
A sequence of polynomials is orthogonal with respect to a weight $w$ when $\int w(x)p_n(x)p_m(x)dx=0$ for $n\ne m$. Legendre, Chebyshev, and Hermite polynomials are standard examples. \par\noindent\textbf{Applications:} orthogonal polynomials support approximation, Gaussian quadrature, spectral PDE methods, and probability.

\term{Orlicz Space}
An Orlicz space generalizes $L^p$ by measuring size through a Young function $\Phi$ rather than only a power. It can describe growth between or beyond standard power laws. \par\noindent\textbf{Applications:} Orlicz spaces are used in nonlinear PDEs, harmonic analysis, elasticity, and models with nonstandard growth.

\term{Orthonormal Set}
An orthonormal set consists of unit vectors that are pairwise orthogonal. Such a set is automatically linearly independent, although it may not span the whole space. \par\noindent\textbf{Applications:} orthonormal sets provide efficient coordinates, signal decompositions, and basis construction.

\term{Orthonormality}
Orthonormality means that vectors are pairwise orthogonal and each has norm one. An orthonormal set becomes an orthonormal basis when it also spans the space. \par\noindent\textbf{Applications:} orthonormality simplifies projections, numerical computation, Fourier analysis, and quantum-state representations.

\term{Outlier}
An outlier is an observation unusually far from the main pattern of a dataset. It may result from error, a rare event, or a genuine subpopulation, so it should not be removed automatically. \par\noindent\textbf{Applications:} outlier analysis supports quality control, fraud detection, robust statistics, and anomaly monitoring.

\term{Outer Product}
The outer product of vectors $u$ and $v$ is the matrix $uv^{\mathsf T}$ with entries $(uv^{\mathsf T})_{ij}=u_iv_j$. Unlike an inner product, it produces a matrix. \par\noindent\textbf{Applications:} outer products create rank-one updates, covariance estimates, tensor decompositions, and attention or feature interactions.

\term{Outer Automorphism}
An outer automorphism is a group automorphism not obtained by conjugation within the group. The quotient $\operatorname{Out}(G)=\operatorname{Aut}(G)/\operatorname{Inn}(G)$ measures automorphisms external to the group's internal symmetries. \par\noindent\textbf{Applications:} outer automorphisms support classification and extension problems in group theory.

\term{Overlap}
An overlap is the part common to two or more sets, intervals, or regions. For measurable sets $A$ and $B$, it is represented by $A\cap B$ and its size by $\mu(A\cap B)$. \par\noindent\textbf{Applications:} overlap measures intersections in geometry, image matching, scheduling, databases, and probability.

\term{Ordinal Number}
An ordinal number describes the order type of a well-ordered set. The first infinite ordinal is $\omega$, the order type of the natural numbers, followed by larger transfinite ordinals. \par\noindent\textbf{Applications:} ordinals support transfinite induction, recursion, and the analysis of well-founded hierarchies in set theory.

\term{Overline}
An overline can denote the closure $\overline A$ of a set, the complex conjugate $\overline z$, or an average depending on context. The closure contains the set and all of its limit points. \par\noindent\textbf{Applications:} overlines describe limits and boundaries in topology, conjugation in complex algebra, and averages in statistics.
